\section{Introduction}\label{sec:intro}

\subsection{Problem}
A bank customer authenticates to the bank, but the device that carries the
customer's intent to the bank --- the web browser on a personal computer --- is
not trustworthy. Malware in the browser, known as a man-in-the-browser (MitB)
Trojan, can read and rewrite every request and response. When the customer asks
to send money to account~A, the Trojan can change the request to account~B while
leaving the page on screen showing account~A. The same manipulation is available
to a real-time phishing proxy that relays the victim's traffic through an
attacker-controlled site. A one-time password (OTP) delivered out of band, such
as a time-based code from an authenticator app, does \emph{not} defend against
this: the code proves that the right person is present, but says nothing about
\emph{which} transaction is being approved, so the attacker simply uses it to
authorise the tampered transfer.

\subsection{Approach}
The defence studied here is \emph{transaction binding}: the code or signature
that the customer returns is computed over the exact transaction, so a code that
was produced for ``send 1{,}000{,}000 to account~A'' cannot authorise ``send
9{,}000{,}000 to account~B''. For this to help a human, the transaction must be
shown to the customer on a second, trusted device (the phone) --- the
``what-you-see-is-what-you-sign'' (WYSIWYS) principle --- so the customer can
notice a mismatch before approving.

We build a demo bank with four interchangeable modes and compare them under one
threat model:
\begin{itemize}
  \item \texttt{baseline} --- a plain TOTP, \emph{not} bound to the transaction;
  \item \texttt{hmac} --- an 8-digit code over the exact transaction (HMAC-SHA256, OCRA-style);
  \item \texttt{sig} --- an ECDSA P-256 signature over the exact transaction;
  \item \texttt{sig\_blind} --- identical cryptography to \texttt{sig}, but the
        phone displays no transaction details.
\end{itemize}
The \texttt{sig\_blind} mode is a deliberate experiment: it keeps the
cryptography of \texttt{sig} but removes the trusted display, so that the
difference between \texttt{sig} and \texttt{sig\_blind} measures the role of the
display alone.

\subsection{Research questions}
\begin{description}
  \item[RQ1.] Does binding the authentication code to the transaction stop MitB
    and real-time phishing better than an ordinary OTP?
  \item[RQ2.] How do symmetric binding (HMAC) and asymmetric binding
    (signatures) differ in security and in cost?
  \item[RQ3.] What does the trusted display (WYSIWYS) contribute, independently
    of the cryptography? The blind-signing mode answers this.
\end{description}

\subsection{Contributions}
\begin{enumerate}
  \item A complete, reproducible demo: a FastAPI bank, a browser-based phone
    authenticator, a shared cryptographic core, and four modes selectable at
    run time (\cref{sec:design,sec:impl}).
  \item A threat model of seven attacker capabilities and seven security goals,
    and an automated harness that checks thirteen scenarios in all four modes
    ($56$ cells) against an expected matrix, so that any regression is caught
    mechanically (\cref{sec:eval}).
  \item A controlled measurement of what each layer buys: transaction binding,
    asymmetric keys, and the trusted display, including a live MitB
    demonstration and micro-benchmarks (\cref{sec:eval,sec:perf}).
  \item An explicit account of the limits: the per-transaction guessing bound,
    the gap between verifiable evidence and non-repudiation, and the idealised
    user model (\cref{sec:discussion}).
\end{enumerate}
We do \emph{not} claim transaction binding as a novel mechanism; it is long
established (\cref{sec:background}). The contribution is the controlled
comparison on one system.
